\section{Implementación y depuración de policy gradient}
\subsection{Pseudocódigo y mejores prácticas}
A continuación se muestra el pseudocódigo típico de policy gradient:
\begin{enumerate}
    \item Iniciar el entorno, inicializar $\theta$, preparar el buffer.
    \item Recolectar $N$ trajectories, registrando $(s_t,a_t,r_t)$.
    \item Calcular el reward-to-go y el baseline (puede usarse una value network).
    \item Construir la surrogate loss y ejecutar una ronda de gradient update.
    \item Repetir, evaluando la política cada $K$ pasos.
\end{enumerate}

\begin{table}[H]
\centering
\small
\begin{tabular}{p{0.32\textwidth}p{0.4\textwidth}p{0.26\textwidth}}
\toprule
Componente & Rango recomendado & Descripción \\
\midrule
Batch size & 2048--8192 steps & Determina la variance y el compute \\
Learning rate & 1e-4 -- 5e-4 & Un learning rate pequeño garantiza estabilidad \\
Entropy coeff & 0.01--0.03 & Incentiva la exploración y evita una convergencia demasiado rápida \\
$\lambda$ & 0.9--0.98 & GAE smoothing \\
\bottomrule
\end{tabular}
\caption{Hiperparámetros clave al entrenar policy gradient}
\end{table}

\subsection{Notas de depuración}
\begin{warningbox}{Ignorar los registros hace perder errores}
El policy gradient es propenso al gradient exploding/vanishing, a la exploding reward scale y al policy collapse. Es indispensable registrar la reward curve, la policy entropy y la KL divergence; de lo contrario resulta difícil ubicar el problema.
\end{warningbox}

\begin{importantbox}{Recomendaciones de depuración}
\begin{itemize}
    \item Verificar primero el pipeline de datos en un entorno pequeño (CartPole);
    \item Monitorear manualmente la distribución del advantage, asegurando que su media se mantenga cerca de cero;
    \item Usar gradient clipping o trust region para evitar actualizaciones con pasos demasiado grandes.
\end{itemize}
\end{importantbox}

\subsection{Resumen de la sección}
A nivel de implementación es necesario encadenar el muestreo, la estimación, la optimización y el registro en un ciclo cerrado; hiperparámetros razonables, métricas de depuración y un mecanismo de registro son la garantía para evitar que el entrenamiento se descomponga.
